% =====================================================================
% 第 VII 章  Construction of class fields
%          (Carl S. Herz, Dec. 18, 1957, revised Nov. 1965)
% =====================================================================
\lecture{Construction of class fields}{Construction of Class Fields}{Carl S. Herz}{Dec. 18, 1957 (revised Nov.\ 1965)}

\setcounter{section}{-1}

% ---------------------------------------------------------------
\section{Introduction}\label{sec:VII-intro}

The purpose of this lecture is to give some explicit constructions of
class fields of imaginary quadratic fields by arithmetic means.

Let $K$ be an algebraic number field and $\mathfrak{H}$ the group of
ideal classes of $K$. The absolute (Hilbert) class field of $K$ is the
class field corresponding to $\mathfrak{H}$; it is the unique maximal
unramified abelian extension $H$ of $K$. In general, it is not easy to
determine the structure of the group $\mathfrak{H}$ and the explicit
construction of the field $H$ is very difficult. The problem is attacked
this way. Suppose $K/k$ is normal with Galois group $\Pi$; then $H/k$ is
normal, and by means of class field theory one can describe the Galois
group $\Gamma$ of $H/k$ exactly in terms of $\Pi$, $\mathfrak{H}$, and
the action of $\Pi$ on $\mathfrak{H}$. Here we shall take $k=\mathbb{Q}$
and restrict our attention to $K/\mathbb{Q}$ cyclic; in this case $\Gamma$
is naturally isomorphic to the semi-direct product $\Pi\mathfrak{H}$. The
rest is algebraic number theory: knowing the Galois group of
$H/\mathbb{Q}$ and that $H/K$ is unramified we can go a long way towards
writing $H$ in the form $H=\mathbb{Q}(\theta_{1},\ldots,\theta_{n})$
where the $\theta_{i}$ are roots of explicitly determined polynomials.

When $K$ is a quadratic field then $\Pi\simeq\mathbb{Z}_{2}$ and the
action of $\Pi$ on $\mathfrak{H}$ is given by
$\mathfrak{k}^{\tau}=\mathfrak{k}^{-1}$,
$\mathfrak{k}\in\mathfrak{H}$ and $\tau$ the generator of $\Pi$. This
situation is very simple indeed. Suppose $L$ is a subfield of $H$ which
is cyclic of degree $\ell$ over $K$; then $L/\mathbb{Q}$ is normal with a
dihedral Galois group. It is easy to describe dihedral extensions of the
rationals; thus the problem is to find such fields $L$ which are
unramified over $K$, and this does not require complete knowledge of the
class group $\mathfrak{H}$. The cases $\ell=2,3,$ or $4$ are simple
enough so that explicit calculations can be carried out.

In summary, if the exponent of the class group $\mathfrak{H}$ of a
quadratic number field $K$ divides $12$, the Hilbert class field $H$ is
easy to construct by arithmetic means. If $K$ is an imaginary quadratic
field we have (V, Theorem 2) $H=K(j(\mathfrak{k}))$ where
$j(\mathfrak{k})$ is any class invariant. Except for very small
discriminants, the calculation of the class invariant is intractable
unless $\mathfrak{H}\simeq\mathbb{Z}_{2}^{\,t-1}$ or
$\mathfrak{H}\simeq\mathbb{Z}_{4}\times\mathbb{Z}_{2}^{\,t-2}$, and in
these cases the arithmetic construction of $H$ is respectively immediate,
almost immediate.

% ---------------------------------------------------------------
\section{Preliminaries from class field theory}\label{sec:VII-prelim}

Let $H$ be the Hilbert class field of $K$.

\begin{lemma}\label{lem:VII-1}
If $K/k$ is normal then $H/k$ is normal.
\end{lemma}

\begin{proof}[Proof.]
Let $\sigma$ be an isomorphism of $H/k$. Then $H^{\sigma}$ is an
unramified abelian extension of $K^{\sigma}=K$. Since $H$ is the unique
maximal abelian unramified extension of $K$, it follows that
$H^{\sigma}\subset H$ and hence $H^{\sigma}=H$.
\end{proof}

Let $\Pi$ denote the Galois group of $K/k$ (action on the right). Then
there is a natural action of $\Pi$ on $\mathfrak{H}$, the group of ideal
classes of $K$.

\begin{namedstmt}{ARTIN RECIPROCITY LAW}
Let $H$ be the class field of $K$ corresponding to the group
$\mathfrak{H}$ of ideal classes. Then the Galois group of $H/K$ is
canonically isomorphic as a $\Pi$-module to $\mathfrak{H}$.
\end{namedstmt}

Let $\Gamma$ denote the Galois group of $H/k$. According to the Artin
Reciprocity Law, $\Gamma$ is an extension of $\mathfrak{H}$ by $\Pi$. The
precise determination of $\Gamma$ depends on the particular element
$\bar{u}\in H^{2}(\Pi,\mathfrak{H})$ describing the extension in question.
According to the Weil--Shafarevitch Theorem, $\bar{u}$ is the image of
$u$, the canonical class of $K/k$, under a natural projection. In order
to avoid going into details here we restrict our attention to
$k=\mathbb{Q}$ (it suffices to assume that all abelian subfields of
$K/k$ are ramified over $k$).

\begin{namedstmt}{COROLLARY OF SHAFAREVITCH--WEIL THEOREM}
If $K/\mathbb{Q}$ is normal with Galois group $\Pi$ then the Galois group
of $H/\mathbb{Q}$ is isomorphic to the semi-direct product
$\Pi\mathfrak{H}$.
\end{namedstmt}

For $K$ an imaginary quadratic field the appeal to the
Shafarevitch--Weil Theorem can be avoided. We have seen in V that
$H=K(j(\mathfrak{k}))=\mathbb{Q}(\sqrt{d},j(\mathfrak{k}))$ where the
class invariant $j(\mathfrak{k})$ is a root of a polynomial with
coefficients in $\mathbb{Q}$ which is irreducible over $K$. Let $\tau$ be
the automorphism of $H/\mathbb{Q}$ given by complex conjugation; the
restriction of $\tau$ to $K$ is the non-trivial automorphism of
$K/\mathbb{Q}$. Moreover, $(j(\mathfrak{k}))^{\tau}
=j(\mathfrak{k})^{\tau}=j(\mathfrak{k}^{\tau})$. Let $\mathfrak{k}_{1}$
be the unit class; then $\tau$ can also be described as the automorphism
of $H/\mathbb{Q}$ determined by
\[
\sqrt{d}\rightarrow-\sqrt{d},\qquad j(\mathfrak{k}_{1})
\rightarrow j(\mathfrak{k}_{1}).
\]
The elements $\mathfrak{k}\in\mathfrak{H}$ are associated with
automorphisms
\[
\sqrt{d}\rightarrow\sqrt{d},\qquad j(\mathfrak{k}_{1})
\rightarrow j(\mathfrak{k}^{-1})
\]
of $H/\mathbb{Q}$. Thus we have $2h$ distinct automorphisms
$\mathfrak{k}_{1},\ldots,\mathfrak{k}_{h}$;
$\tau\mathfrak{k}_{1},\ldots,\tau\mathfrak{k}_{h}$ of $H/\mathbb{Q}$
where $h=\#\mathfrak{H}$ = class number of $K$. Since
$[H:\mathbb{Q}]=2h$, these automorphisms exhaust $\Gamma$, the Galois
group of $H/\mathbb{Q}$. By direct calculation we have
$\sqrt{d}{}^{\,\tau^{-1}\mathfrak{k}\tau}=\sqrt{d}$ and
\[
j(\mathfrak{k}_{1})^{\tau^{-1}\mathfrak{k}\tau}
=j(\mathfrak{k}_{1})^{\mathfrak{k}\tau}
=(j(\mathfrak{k}^{-1}))^{\tau}
=j(\mathfrak{k}^{-\tau}),
\]
which proves that $\tau^{-1}\mathfrak{k}\tau=\mathfrak{k}^{\tau}$. Hence
$\Gamma$ is isomorphic to the semi-direct product $\Pi\mathfrak{H}$.

% ---------------------------------------------------------------
\section{The genus field and the group of genera}\label{sec:VII-genus}

Let $K/k$ be abelian with Galois group $\Pi$, and let $H$ be the Hilbert
class field of $K$. The \emph{genus field} of $K/k$ is the maximal
abelian subfield $G$ of $H/k$. If $\Gamma$ is the Galois group of $H/k$
then $G$ is the fixed field of $\Gamma'$, the commutator subgroup of
$\Gamma$. Since $K/k$ is abelian, the canonical image of $\mathfrak{H}$
in $\Gamma$ must contain $\Gamma'$. Thus $\Gamma'$ is canonically
isomorphic as a $\Pi$-module with a subgroup $\mathfrak{G}$ of
$\mathfrak{H}$; this subgroup is called the principal genus.

\begin{lemma}\label{lem:VII-2}
The principal genus $\mathfrak{G}$ is the subgroup of $\mathfrak{H}$
generated by the elements $\mathfrak{k}^{1-\tau}$ where
$\mathfrak{k}\in\mathfrak{H}$, $\tau\in\Pi$.
\end{lemma}

\begin{proof}[Proof.]
We know that $\Gamma$ is canonically isomorphic to an extension of
$\mathfrak{H}$ by $\Pi$. Let $u_{\tau}\in\Gamma$ be a representative of
$\tau\in\Pi$. The elements of $\Gamma$ are identified with the products
$u_{\tau}\mathfrak{k}$ and
$u_{\tau}^{-1}\mathfrak{k}^{-1}u_{\tau}\mathfrak{k}
=\mathfrak{k}^{1-\tau}$. These commutators generate $\Gamma'$ since
$\mathfrak{H}$ is abelian.
\end{proof}

The factor group $\mathfrak{H}_{\Pi}=\mathfrak{H}/\mathfrak{G}$ is called
the group of genera of $K/k$. It is canonically isomorphic under the
reciprocity map with the Galois group of $H/G$.

The essential fact used in the construction of the genus field is given by

\begin{namedstmt}{PROPOSITION 1}
Let $K/k$ be abelian. The genus field of $K/k$ is the unique maximal
abelian extension $G/k$ with the same ramification as $K/k$.
\end{namedstmt}

\begin{proof}[Proof.]
Let $I_{k}$ denote the group of ideals of $k$ relatively prime to the
discriminant of $K/k$ and $P_{K}$ the group of principal ideals of $K$,
again relatively prime to the discriminant of $K/k$, an assumption made
for all ideals in the course of this proof. An abelian extension $L/k$
has the same ramification as $K/k$ if
$N_{K/k}P_{K}=N_{L/k}P_{L}$ as subgroups of $I_{k}$. Now let $H$ be the
Hilbert class field to $K$. By definition, $N_{H/K}I_{H}=P_{K}$, but by
the Principal Ideal Theorem, $N_{H/K}P_{H}=P_{K}$. Since the genus field
$G$ is the maximal abelian subfield of $H/k$ we have
\[
N_{K/k}\,P_{K}=N_{K/k}\,N_{H/K}\,P_{H}=N_{H/k}\,P_{H}
=N_{G/k}\,N_{H/G}\,P_{H}\subset N_{G/k}\,P_{G},
\]
\[
N_{K/k}\,P_{K}=N_{K/k}\,N_{H/K}\,I_{H}=N_{H/k}\,I_{H}
=N_{G/k}\,N_{H/G}\,I_{H}=N_{G/k}\,I_{G}.
\]
and so $N_{K/k}P_{K}=N_{G/k}I_{G}=N_{G/k}P_{G}$. Thus $G/k$ has the same
ramification as $K/k$, and $G/k$ is the class field to the subgroup
$N_{K/k}P_{K}$ of $I_{k}$. Since for any extension $L/k$ we have
$N_{L/k}P_{L}\subset N_{L/k}I_{L}$, it follows that for any abelian
extension $L/k$ with the same ramification as $K/k$ we have
$N_{G/k}I_{G}\subset N_{L/k}I_{L}$; so $L$ is a subfield of $G$.
\end{proof}

Proposition~1 leads to a simple effective procedure for constructing
genus fields. To do this we examine $N_{K/k}P_{K}$. Given a prime $p$ of
$k$ let $U_{p}$ denote the group of elements of $k$ relatively prime to
$p$ and let $V_{p}$ be the subgroup of $U_{p}$ consisting of those
elements which are congruent to norms from $K$ modulo arbitrarily high
powers of $p$. Then a necessary and sufficient condition that
$(\alpha)\in N_{K/k}P_{K}$ is
\[
\alpha\in E_{0}\cdot\bigcap V_{p}
\]
where $E_{0}$ is the group of units of $k$ and the intersection is taken
over primes $p$ dividing the discriminant of $K/k$, including possibly
real infinite primes where ``congruent'' is to be understood as ``has the
same sign''. In case $K/k$ is cyclic of prime degree $n$ it follows that
\[
U_{p}^{\,n}\subset V_{p}\qquad\text{and}\qquad[U_{p}:V_{p}]=n,
\]
where $U_{p}^{\,n}$ is the subgroup of $U_{p}$ consisting of elements
which are congruent to n-th powers modulo arbitrarily high powers of $p$.

The situation is particularly simple when $K$ is a cyclic extension of
$p\equiv 0,1\bmod n$ since $U_{p}^{\,n}=U_{p}$ unless $n\mid\varphi(p^{r})$
where $\varphi(p^{r})$ is the order of $\mathbb{Z}_{p^{r}}^{*}$, the
group of units modulo $p^{r}$. Moreover, $\mathbb{Z}_{p^{r}}^{*}$ is
cyclic for $p\neq 2$. Hence
\[
[U_{p}:U_{p}^{\,n}]=n\qquad\text{for }p\equiv 0,1\bmod n,
\]
with the single exception
\[
U_{2}/U_{2}^{2}\simeq\mathbb{Z}_{2}^{2}.
\]
At $\infty$ there is no ramification when $n$ is odd, for then
$U_{\infty}^{\,n}=U_{\infty}$; but $[U_{\infty}:U_{\infty}^{2}]=2$.

We conclude that if $K/\mathbb{Q}$ is cyclic of prime degree $n$ and $p$
ramifies then $V_{p}=U_{p}^{\,n}$ for all $p$ finite or infinite except
for $p=2$, which can only occur when $n=2$. There are three possibilities
in the exceptional case:
\[
V_{2}=\{1,v\}\cdot U_{2}^{2}\quad\text{where }v=-1,\ 5,\text{ or }3.
\]

It is always true that $P_{\mathbb{Q}}/N_{K/\mathbb{Q}}P_{K}
\simeq\bigcap U_{p}/E_{0}\cap V_{p}$ where $p$ runs over the ramified
primes, or, for that matter, any finite set of primes containing all the
ramified primes. Using the Chinese Remainder Theorem, we obtain a direct
product decomposition best written in the form (all isomorphisms being
canonical):
\[
P_{\mathbb{Q}}/N_{K/\mathbb{Q}}P_{K}
\simeq\Bigl(\bigcap U_{p}\Bigr)\big/\bigl(E_{0}\cap V_{p}\bigr)
\simeq\prod\Gamma_{p},
\]
where $\Gamma_{p}=U_{p}/V_{p}$ if $E_{0}\subset V_{p}$ and
$\Gamma_{p}=(U_{p_{0}}\cap U_{p})/E_{0}(V_{p_{0}}\cap V_{p})$ if
$E_{0}\notin V_{p}$, $p_{0}$ being a fixed ramified prime such that
$-1\notin V_{p_{0}}$. The reason for writing the product this way is that
there is a unique extension $\mathbb{Q}(\theta_{p})/\mathbb{Q}$ which is
cyclic of degree $n$ and whose ramification is described by $\Gamma_{p}$,
namely the class field of $\mathbb{Q}$ corresponding to $\Gamma_{p}$.
Since the genus field $G$ is the class field of $\mathbb{Q}$
corresponding to $\prod\Gamma_{p}$, it must be the compositum of the
$\mathbb{Q}(\theta_{p})$. What we have proved is

\begin{namedstmt}{THEOREM 1}
Let $K/\mathbb{Q}$ be cyclic of prime degree $n$. The genus field of
$K/\mathbb{Q}$ is $G=K(\theta_{1},\ldots,\theta_{t})$ where
$p_{1},\ldots,p_{t}$ are the distinct prime divisors of the discriminant
$d$ (excluding $p_{0}=\infty$ if $K$ is imaginary quadratic and an
arbitrary $p_{0}\equiv -1\bmod 4$ if $K$ is real quadratic) where for
each $p_{i}$, $\mathbb{Q}(\theta_{i})$ is cyclic of degree $n$ over
$\mathbb{Q}$ and specified this way:

(i) for $n$ odd and $p$ given (so $p\equiv 0,1\bmod n$), $\mathbb{Q}
(\theta)$ is the unique cyclic extension of degree $n$ which ramifies
only at $p$;

(ii) for $n=2$ and $p\equiv 1\bmod 4$, $\theta=\sqrt{p}$;

(iii) for $n=2$ and $p\equiv -1\bmod 4$, $\theta=\sqrt{-p}$ or
$\theta=\sqrt{p_{0}p}$ according to the imaginary or real case;

(iv) for $n=2$ and $p=2$, so $K=\mathbb{Q}(\sqrt{d})$, the cases are
\[
\begin{aligned}
&\theta=\sqrt{2}\ \text{if } d\equiv 8\bmod 32\\
&\theta=\sqrt{-1}\ \text{or}\ \sqrt{p_{0}}\ \text{if } d\equiv 4\bmod 8
\ \text{according to } d<0 \text{ or } d>0\\
&\theta=\sqrt{-2}\ \text{or}\ \sqrt{2p_{0}}\ \text{if }
d\equiv -8\bmod 32\ \text{according to } d<0 \text{ or } d>0.
\end{aligned}
\]
\end{namedstmt}

It may not be obvious how we arrived at (iv). To see it one observes that
if $2$ ramifies then $K=\mathbb{Q}(\sqrt{d_{0}})$ where the discriminant
is $4d_{0}$. Then $1-d_{0}\in V_{2}$ since $1-d_{0}$ is the norm of
$1+\sqrt{d_{0}}$. If $d_{0}\equiv 2\bmod 8$ then $-1\in V_{2}$ and
$\Gamma_{2}=U_{2}/V_{2}$ corresponds to $\mathbb{Q}(\sqrt{2})$. The
remaining possibilities are $d_{0}\equiv -2\bmod 8$ and
$d_{0}\equiv -1\bmod 4$. In the former case $3\equiv 1-d_{0}\bmod 8$, so
$3\in V_{2}$; in the latter $5\equiv 1-4d_{0}\bmod 8$, so $5\in V_{2}$.
In neither of these cases do we have $-1\in V_{2}$; hence
$\Gamma_{2}=U_{p_{0}}\cap U_{2}/E_{0}(V_{p_{0}}\cap V_{2})$. If $d<0$ we
can take $p_{0}=\infty$ and the class field corresponding to $\Gamma_{2}$
is $\mathbb{Q}(\sqrt{-1})$ or $\mathbb{Q}(\sqrt{-2})$ according to
whether $5\in V_{2}$ or $3\in V_{2}$. When $d>0$ we must use
$\mathbb{Q}(\sqrt{p_{0}})$ and $\mathbb{Q}(\sqrt{2p_{0}})$ instead where
$p_{0}\equiv -1\bmod 4$.

\noindent\emph{Remark.} In the classical treatment of real quadratic
fields $K$ one considers $H^{+}$, the maximal abelian extension of $K$
which is unramified except at $\infty$. One puts $G^{+}$ the maximal
abelian subfield of $H^{+}/\mathbb{Q}$ and then
$G^{+}=\mathbb{Q}(\sqrt{p_{1}},\ldots,\sqrt{p_{s}})$ where
$p_{1},\ldots,p_{s}$ are the distinct prime divisors of the discriminant.
Therefore either $G^{+}=G$ or $G^{+}=G(\sqrt{p_{0}})$ according to
whether or not there is a prime $p_{0}\equiv -1\bmod 4$ dividing the
discriminant.

\medskip

The Galois group of $G/\mathbb{Q}$ is $\Gamma/\Gamma'$ and, as we have
just seen, this is isomorphic to $\prod_{p\mid d}\Gamma_{p}$. The
reciprocity law makes this isomorphism explicit. Assume that the action
of $\Pi$ has been extended to $G$. Then $\Pi$ acts on each
$\mathbb{Q}(\theta_{p})$; we shall write $p=p'$ if $\Pi$ acts
non-trivially, $p=p''$ if $\Pi$ acts trivially. Let $q$ be a prime such
that $q\nmid d$, $q\notin V_{p'}$, $q\in V_{p''}$ (assume $q>0$ and
$q\in V_{p_{0}}$ in the quadratic cases). Then $q$ determines a generator
$\gamma_{p'}$ of $\Gamma_{p'}$. The reciprocity law for the extension
$\mathbb{Q}(\theta_{p'})/\mathbb{Q}$ says
$\gamma_{p'}\Bigl(\frac{\mathbb{Q}(\theta_{p'})/\mathbb{Q}}{q},q\Bigr)=1$.
The interpretation of the reciprocity symbol gives the fact that $q$
remains prime in $\mathbb{Q}(\theta_{p'})/\mathbb{Q}$ and $\gamma_{p'}$
induces the Frobenius automorphism on the residue class extension at
$q$. This determines $\gamma_{p'}$ uniquely. It follows that
$\Bigl(\frac{G/\mathbb{Q}}{q},q\Bigr)=(\prod\gamma_{p'})^{-1}$. On the
other hand $q$ must remain prime in $K/\mathbb{Q}$ since its divisors in
$G/\mathbb{Q}$ are fixed under $\Pi$. Hence
$\Bigl(\frac{K/\mathbb{Q}}{q},q\Bigr)^{-1}$ is the generator of $\Pi$
which induces the Frobenius automorphism on the residue class extension
at $q$. Clearly $\Bigl(\frac{K/\mathbb{Q}}{q},q\Bigr)$ is the image of
$\Bigl(\frac{G/\mathbb{Q}}{q},q\Bigr)$ under the natural projection of
the Galois group of $G/\mathbb{Q}$ onto the Galois group of
$K/\mathbb{Q}$. Hence
$\tau=\prod_{p'\mid d}\gamma_{p'}$ is a generator of $\Pi$. If the action
of $\Pi$ on $G$ has not been given in advance it suffices to choose $q$
such that $q$ remains prime in $K/\mathbb{Q}$ and determine the $p'$ as
those $p\mid d$ such that $q\notin V_{p}$.

To summarize, we state

\begin{namedstmt}{THEOREM 2}
Let $K/\mathbb{Q}$ be cyclic of prime degree $n$ with Galois group $\Pi$.
Given an action of $\Pi$ on the genus field
$G=\mathbb{Q}(\theta_{1},\ldots,\theta_{t})$ and the reciprocity law
identification of $\Gamma_{p_{i}}$ with the Galois group of
$\mathbb{Q}(\theta_{i})/\mathbb{Q}$, there is a choice of generators
$\gamma_{p}$ for $\Gamma_{p}$ such that $\tau=\prod_{p'}\gamma_{p'}$
generates $\Pi$ and the group of genera is
$\mathfrak{H}_{\Pi}=(\prod_{p\mid d}\Gamma_{p})/\Pi$, where $p'$
denotes those $p_{i}$ such that $\Pi$ acts non-trivially on $\theta_{i}$.
\end{namedstmt}

For $K$ imaginary quadratic one takes $\tau$ to be complex conjugation on
$H$, hence also on $G$. Then the $p'$ are the $p\mid d$ such that
$p\equiv -1\bmod 4$ and $2$ if $d\equiv 4\bmod 8$ or $d\equiv -8\bmod 32$.

The concrete isomorphism of $\mathfrak{H}_{\Pi}$ with
$(\prod_{p\mid d}\Gamma_{p})/\Pi$ is made this way. Each genus can be
represented by a prime $\mathfrak{q}$ of $K$, but since conjugates belong
to the same genus it suffices to specify the rational prime $q$ which
$\mathfrak{q}$ divides. The groups in question are isomorphic to
$\mathbb{Z}_{n}^{\,t-1}$; so it remains only to give rational primes
$q_{1},\ldots,q_{t-1}$ representing generators of $\mathfrak{H}_{\Pi}$ in
correspondence with certain elements of $\prod_{p\mid d}\Gamma_{p}$.
Given $\gamma_{p}$ as an automorphism of $G/\mathbb{Q}$, let
$\bar{\gamma}_{p}$ be the automorphism of $K/\mathbb{Q}$ obtained from the
natural projection of the Galois group of $G/\mathbb{Q}$ onto the Galois
group of $K/\mathbb{Q}$. Take $p_{t}$ so that $\bar{\gamma}_{p_{t}}$ is
non-trivial (this is no restriction at all). Since $\Pi$ is cyclic of
prime degree we can find integers $e_{1},\ldots,e_{t-1}$ such that
$\bar{\gamma}_{p_{i}}\bar{\gamma}_{p_{t}}^{e_{i}}=1$. Now consider the
$\gamma_{p}$ as elements of the groups $\Gamma_{p}$ defined by certain
congruences in $\mathbb{Q}$ and choose primes $q_{1},\ldots,q_{t-1}$
($q_{i}>0$ and $q_{i}\in V_{p_{0}}$ in the quadratic cases) such that
\[
q_{i}\in\gamma_{p_{i}}V_{p_{i}}\cap\gamma_{p_{t}}^{e_{p_{i}}}V_{p_{t}}
\cap\bigcap_{j\neq i,t}V_{p_{j}}.
\]

\begin{namedstmt}{THEOREM 3}
The canonical isomorphism $(\prod_{p\mid d}\Gamma_{p})/\Pi
\to\mathfrak{H}_{\Pi}$ is completely specified by
$\gamma_{p_{i}}\gamma_{p_{t}}^{e_{i}}\to q_{i}$, $i=1,\ldots,t-1$. Here
$(\prod_{p\mid d}\Gamma_{p})/\Pi$ is to be viewed as the Galois group of
$G/K$ and $\mathfrak{H}_{\Pi}$ as the group of genera of $K$.
\end{namedstmt}

\begin{proof}[Proof.]
The $\gamma_{p_{i}}\gamma_{p_{t}}^{e_{i}}$ generate a subgroup $\Delta$
of $\prod_{p\mid d}\Gamma_{p}$. For $\delta\in\Delta$ we have
$\bar{\delta}=1$, but $\bar{\tau}\neq 1$ where $\tau$ generates $\Pi$.
Thus $\Delta\cap\Pi=\{1\}$, and
$\Delta\simeq\mathbb{Z}_{n}^{\,t-1}\simeq
(\prod_{p\mid d}\Gamma_{p})/\Pi$. Therefore we may identify $\Delta$ with
the Galois group of $G/K=K(\theta_{1},\ldots,\theta_{t-1})/K$.
According to the reciprocity law, $q_{i}$ splits in $K/\mathbb{Q}$ and in
each $\mathbb{Q}(\theta_{j})/\mathbb{Q}$ for $j\neq i,t$ but does not
split in $\mathbb{Q}(\theta_{i})/\mathbb{Q}$. It follows that $q_{i}$
splits completely in
$K(\theta_{1},\ldots,\theta_{i-1},\theta_{i+1},\ldots,\theta_{t-1})/K$
and remains prime in $K(\theta_{i})/K$, where $\mathfrak{q}_{i}$ is any
prime of $K$ dividing $q_{i}$. The reciprocity law demands that we
associate $q_{i}$ with the automorphism of
$K(\theta_{1},\ldots,\theta_{t-1})/K$ which leaves $\theta_{j}$ fixed for
$j\neq i$ and which induces the Frobenius automorphism in the
residue-class extension at $\mathfrak{q}_{i}$ of $K(\theta_{i})/K$. Since
$q_{i}$ splits in $K/\mathbb{Q}$, the residue-class extension at
$\mathfrak{q}_{i}$ of $K(\theta_{i})/K$ is the same as the residue-class
extension at $q_{i}$ of $\mathbb{Q}(\theta_{i})/\mathbb{Q}$. The
automorphism $\gamma_{i}$ of $\mathbb{Q}(\theta_{i})/\mathbb{Q}$
corresponds to $\gamma_{i}\gamma_{t}^{e_{i}}\in\Delta$ when $\Delta$ is
taken as the Galois group of
$K(\theta_{1},\ldots,\theta_{t-1})/K$.
\end{proof}

The above three theorems give a complete description of the theory of
genera for $K/\mathbb{Q}$ cyclic of prime degree $n$. This is equivalent
to the theory of the representation of integers by certain homogeneous
forms of degree $n$ in $n$ rational variables with integral coefficients.
The explicit computations are tedious when $n>2$. In the quadratic case
the computations are nearly trivial, especially since there is no choice
of generators to worry about.

For contrast, we shall give one illustration of Theorems 1, 2, and 3 in
the cubic case.

\medskip\noindent
\emph{Illustration.} Take $K=\mathbb{Q}(\theta)$ where
$\theta^{3}-21\theta-28=0$. Ramification occurs at $p_{1}=3$,
$p_{2}=7$. The genus field is, according to Theorem~1,
$G=\mathbb{Q}(\theta_{1},\theta_{2})$ where $\theta_{1}^{3}-3\theta_{1}-1=0$
and $\theta_{2}^{3}-3\cdot 7\,\theta_{2}-7=0$. We see that $q=5$ remains
prime in $\mathbb{Q}(\theta)$, $\mathbb{Q}(\theta_{1})$, and
$\mathbb{Q}(\theta_{2})$. Hence, by Theorem~2, $\tau=\gamma_{3}\gamma_{7}$
where $\gamma_{3}$ corresponds to elements $\equiv\pm 4\bmod 9$ and
$\gamma_{7}$ to elements $\equiv\pm 2\bmod 7$. It necessarily follows
that $\bar{\gamma}_{3}=\tau^{2}$, $\bar{\gamma}_{7}=\tau^{2}$. According
to Theorem~3 the non-principal genera are given by
$59\leftrightarrow\gamma_{3}\gamma_{7}^{2}$ and
$2\leftrightarrow\gamma_{3}^{2}\gamma_{7}$. We have
$\mathfrak{H}_{\Pi}\simeq\mathbb{Z}_{3}$. (Also
$\mathfrak{H}=\mathfrak{H}_{\Pi}$.) The theory states that $m\in U_{3}\cap U_{7}$ is
the norm of a principal ideal, i.e.\ there exist $x,y,z\in\mathbb{Q}$
such that
\[
m=x^{3}+7\bigl\{4y^{3}+16\cdot 7z^{3}+6x^{2}z-3xy^{2}
+9\cdot 7xz^{2}-12\cdot 7yz^{2}\bigr\},
\]
iff $m\equiv\pm 1\bmod 9$ and $m\equiv\pm 1\bmod 7$. The cubic
reciprocity law states that for a prime $q\neq 3,7$ the congruence
\[
28\equiv w^{3}-21w\bmod q
\]
has a solution iff $q$ corresponds to $\gamma_{3}^{a}\gamma_{7}^{b}$
where $\bar{\gamma}_{3}^{a}\bar{\gamma}_{7}^{b}=1$, i.e.\
$q\equiv\pm 2^{k}\bmod 63$. This is equivalent to $q$ being the norm of
an ideal. Thus $2$ and $59$ are norms of non-principal ideals while
$2\cdot 59$ is the norm of a principal ideal. The presentation of $\tau$,
$\gamma_{3}$, and $\gamma_{7}$ as automorphisms of the appropriate fields
is too tedious.

\medskip

Let us return to the quadratic case. Let $d$ be the discriminant of a
quadratic field $K=\mathbb{Q}(\sqrt{d})$ and $m$ a positive integer
relatively prime to $d$ (also assume $\bigl(\frac{m}{p_{0}}\bigr)=1$).
We have seen that $m$ is the norm of a principal ideal in $K$, i.e.
\[
x^{2}-dy^{2}=4m
\]
has solutions $x,y\in\mathbb{Q}$ iff $\bigl(\frac{m}{p}\bigr)=1$ for each
odd $p\mid d$ and $m\equiv 1,v\bmod 8$ if $2\mid d$. There is one class
in each genus, i.e.\ $\mathfrak{H}=\mathfrak{H}_{\Pi}$, iff whenever
solutions $x,y\in\mathbb{Q}$ exist then solutions $x,y\in\mathbb{Z}$
exist. Here $\bigl(\frac{m}{p}\bigr)=+1$ if $m$ is a quadratic residue
mod $p$; and if $m$ is a prime, the splitting of $m$ in $K/\mathbb{Q}$ is
determined by the classical quadratic reciprocity law.

\medskip\noindent
\emph{Example 1.} Take $d=-84$. Then $H=G=\mathbb{Q}(\sqrt{-1},
\sqrt{-3},\sqrt{-7})$. We have $\tau=\gamma_{2}\gamma_{3}\gamma_{7}$.
Generators for $\mathfrak{H}_{\Pi}$ are given by
\[
19\leftrightarrow\gamma_{2}\gamma_{7}\qquad
5\leftrightarrow\gamma_{3}\gamma_{7}.
\]
The equation $x^{2}+21y^{2}=m$ has integral solutions whenever it has
rational solutions.

\medskip\noindent
\emph{Example 2.} Take $d=-87$. Then
$G=\mathbb{Q}(\sqrt{-3},\sqrt{29})$. We have $\tau=\gamma_{3}$ and a
generator for $\mathfrak{H}_{\Pi}$ is given by
$2\leftrightarrow\gamma_{3}\gamma_{29}$. Notice that $7$ corresponds to
the principal genus, and, indeed, the equation $x^{2}+87y^{2}=4\cdot 7$
has the solution $x=5/2$, $y=1/2$ but it does not have integral
solutions. Hence the primes dividing $7$ give non-principal classes in
the principal genus. On the other hand $16^{2}+87\cdot 1^{2}=7^{3}$.
Therefore if $\mathfrak{P}_{7}$ is a prime of $K$ dividing $7$ then
$\mathfrak{P}_{7}$ is in the principal genus and $\mathfrak{P}_{7}^{3}$
is in the principal class. In fact
$\mathfrak{H}\simeq\mathbb{Z}_{2}\times\mathbb{Z}_{3}$. We shall see
later that $H=\mathbb{Q}(\sqrt{-3},\sqrt{29},\varepsilon^{1/3})$ where
$\varepsilon^{2}-5\varepsilon-1=0$.

\medskip\noindent
\emph{Example 3.} Take $d=-39$. Then $G=\mathbb{Q}(\sqrt{-3},
\sqrt{13})$. We have $\tau=\gamma_{3}$ and a generator for
$\mathfrak{H}_{\Pi}$ is given by $2\leftrightarrow\gamma_{3}\gamma_{13}$.
Let $\mathfrak{P}_{2}$ be a prime of $K$ dividing $2$. Obviously
$\mathfrak{P}_{2}^{2}$ is in the principal genus, but $\mathfrak{P}_{2}^{2}$
is not principal. To see this, observe that
$\mathfrak{P}_{2}\neq\mathfrak{P}_{2}^{\tau}$ since $2$ is unramified.
Hence $\mathfrak{P}_{2}^{2}\neq(2)=\mathfrak{P}_{2}^{1+\tau}$; so if
$\mathfrak{P}_{2}^{2}$ were principal there would be a solution to
$x^{2}+39y^{2}=4\cdot 2^{2}$ with $y\neq 0$, $y\in\mathbb{Z}$. Here
$\mathfrak{H}\simeq\mathbb{Z}_{4}$.

\medskip\noindent
\emph{Example 4.} Take $d=40$. Here $G=\mathbb{Q}(\sqrt{2},\sqrt{5})$.
We have $\tau=\gamma_{5}$. The generator for $\mathfrak{H}_{\Pi}$ is
given by $3\leftrightarrow\gamma_{2}\gamma_{5}$. The equation
$x^{2}-10y^{2}=m$ has integral solutions whenever it has rational
solutions and there is no need to take $m>0$ since
$3^{2}-10\cdot 1^{2}=-1$.

\medskip\noindent
\emph{Example 5.} Take $d=21$. Here $G=K$. The equation
$x^{2}-21y^{2}=4m$ has integral solutions for all $m$ such that
$\bigl(\frac{m}{3}\bigr)=1$ and $\bigl(\frac{m}{7}\bigr)=1$. Therefore
$x^{2}-21y^{2}=\pm 4m$ has solutions whenever
$\bigl(\frac{m}{3}\bigr)\bigl(\frac{m}{7}\bigr)=1$. If we impose the
restriction $m>0$ there is a change: the ideal $(17)$ is the norm of a
principal ideal but $17$ is not a norm. In this example
$\mathfrak{H}=\mathfrak{G}$, but
$H^{+}=\mathbb{Q}(\sqrt{-3},\sqrt{-7})$.

\medskip\noindent
\emph{Example 6.} Take $d=136$. Here $G=\mathbb{Q}(\sqrt{2},\sqrt{17})$.
There is a choice of $\tau=\gamma_{2}$ or $\tau=\gamma_{17}$. The
equation $x^{2}-34y^{2}=m$ has rational solutions for all
$m\equiv\pm 1\bmod 8$, $\equiv\pm 2^{k}\bmod 17$. Whenever
$x^{2}-34y^{2}=m$ has rational solutions then $x^{2}-34y^{2}=\pm m$ has
integral solutions but not necessarily with the $+$ sign. The simplest
example is $(5/3)^{2}-34(1/3)^{2}=-1$; there is no integral solution to
$x^{2}-34y^{2}=-1$. In this example $\mathfrak{H}^{+}\simeq\mathbb{Z}_{4}$
while $\mathfrak{H}^{+}\!_{\Pi}=\mathfrak{H}_{\Pi}
=\mathfrak{H}\simeq\mathbb{Z}_{2}$.

% ---------------------------------------------------------------
\section{Structure of the group of ideal classes}\label{sec:VII-structure}

Let $K/\mathbb{Q}$ be an abelian extension of order $n$. Recall that
$\mathfrak{H}$ denotes the group of ideal classes of $K$ and
$\mathfrak{G}$ the principal genus, i.e.\ the subgroup of $\mathfrak{H}$
generated by the $\mathfrak{k}^{1-\tau}$, $\tau\in\Pi$ where $\Pi$ is the
Galois group of $K/\mathbb{Q}$. Put $N=\sum_{\tau\in\Pi}\tau$. Then for
any $\mathfrak{k}\in\mathfrak{H}$ we have that
$\mathfrak{k}^{n-N}\in\mathfrak{G}$. On the other hand,
$\mathfrak{k}^{N}$ corresponds to the ideal class of the norm of an ideal
in $\mathfrak{k}$ and this, being an ideal of $\mathbb{Q}$, is principal.
It follows that $\mathfrak{H}^{n}$, the group of n-th powers in
$\mathfrak{H}$, is a subgroup of $\mathfrak{G}$. Hence the group of
genera $\mathfrak{H}_{\Pi}=\mathfrak{H}/\mathfrak{G}$ is a subgroup of
$\mathfrak{H}/\mathfrak{H}^{n}$.

In case $n$ is a prime we know $\mathfrak{H}_{\Pi}$ as the Galois group
of $G/K$ where the genus field $G$ is given by Theorem~1. The fundamental
numerical invariant here was $t$ where

\smallskip
\centerline{$t=$ number of distinct prime factors of the discriminant
$d$ of $K/\mathbb{Q}$}

\smallskip
\centerline{(excluding one prime factor $p_{0}\equiv -1\bmod 4$ if $K$
is real quadratic).}
\smallskip

\begin{namedstmt}{THEOREM 4}
Let $K/\mathbb{Q}$ be cyclic of prime degree $n$. Then
$\mathfrak{H}_{\Pi}\simeq\mathbb{Z}_{n}^{\,t-1}$. For $n=2$,
$\mathfrak{H}_{\Pi}\simeq\mathfrak{H}/\mathfrak{H}^{2}$; for $n>2$,
$\mathfrak{H}_{\Pi}$ is a quotient group of $\mathfrak{H}/\mathfrak{H}^{n}$
which is non-trivial if $\mathfrak{H}^{n}\neq\mathfrak{H}$.
\end{namedstmt}

\begin{namedstmt}{COROLLARY}
Let $h$ be the class number of $K$. Then $n\mid h$ iff $t>1$.
\end{namedstmt}

\begin{proof}[Proof of Theorem 4.]
According to Theorem~1 the Galois group of $G/\mathbb{Q}$ is the direct
product of $t$ cyclic groups of order $n$. This shows that
$\mathfrak{H}_{\Pi}\simeq\mathbb{Z}_{n}^{\,t-1}$. When $n=2$,
$N=1+\tau$. Hence $\mathfrak{k}^{1-\tau}=\mathfrak{k}^{2}$, since
$\mathfrak{k}^{\tau}=\mathfrak{k}^{-1}$. Therefore
$\mathfrak{G}=\mathfrak{H}^{2}$ for $n=2$. When $n>2$ we cannot conclude
that $\mathfrak{G}=\mathfrak{H}^{n}$ but $\Pi$ acts on $\mathfrak{H}$
leaving $\mathfrak{H}^{n}$ invariant. Thus the cyclic group of order
$n$, $\Pi$, acts on the n-group $\mathfrak{H}/\mathfrak{H}^{n}$. If
$\mathfrak{H}\neq\mathfrak{H}^{n}$ there must be a non-trivial fixed
point of $\Pi$ in $\mathfrak{H}/\mathfrak{H}^{n}$; this implies that
$\mathfrak{G}\neq\mathfrak{H}$.
\end{proof}

To elaborate the last statement further, let $\mathfrak{H}^{\Pi}$ denote
the subgroup of $\mathfrak{H}$ composed of ambiguous ideal classes,
ideal classes $\mathfrak{k}$ which are invariant under $\Pi$, i.e.\
$\mathfrak{k}^{\tau}=\mathfrak{k}$. The sequence
\[
1\longrightarrow\mathfrak{H}^{\Pi}\longrightarrow\mathfrak{H}
\xrightarrow{\ 1-\tau\ }\mathfrak{H}\longrightarrow\mathfrak{H}_{\Pi}
\longrightarrow 1
\]
is exact; hence $\mathfrak{H}^{\Pi}$ and $\mathfrak{H}_{\Pi}$ have the
same order (observe that $\mathfrak{G}\simeq
\mathfrak{H}/\mathfrak{H}^{\Pi}$). For
$\mathfrak{k}\in\mathfrak{H}^{\Pi}$, however,
$\mathfrak{k}^{n}=\mathfrak{k}^{N}=1$. Since $n$ is a prime, it follows
that $\mathfrak{H}^{\Pi}\simeq\mathfrak{H}_{\Pi}$. One may argue directly
that if $\mathfrak{H}$ contains an element of order $n$ then
$\mathfrak{H}^{\Pi}$, the group of fixed points, must be non-trivial.

\noindent\emph{Remark.} The old-fashioned treatment of genera of
quadratic fields was based on the isomorphism of $\mathfrak{H}^{\Pi}$
with $\mathfrak{H}_{\Pi}$. One proved directly that
$\mathfrak{H}^{\Pi}\simeq\mathbb{Z}_{n}^{\,t-1}$. For imaginary quadratic
fields the analysis of $\mathfrak{H}^{\Pi}$ is quite elementary and
yields the construction of the genus field $G$ in a very straightforward
fashion. It gives no information, however, about the representation of
$\mathfrak{H}_{\Pi}$ as automorphisms of $G/K$. In the real quadratic
case, the analysis of $\mathfrak{H}^{\Pi}$ involves consideration of the
units of $K$, but this is ultimately irrelevant.

\medskip

For $n>2$ the action of $\Pi$ on $\mathfrak{H}$ may be complicated. For
$n=2$ it is just $\mathfrak{k}^{\tau}=\mathfrak{k}^{-1}$. Therefore we
have immediately

\begin{namedstmt}{THEOREM 5}
Let $K/\mathbb{Q}$ be a quadratic field and $L$ an unramified abelian
extension of $K$. Then $L$ is normal over $\mathbb{Q}$. Moreover if
$L/K$ is cyclic then $L/\mathbb{Q}$ is dihedral.
\end{namedstmt}

A consequence of Theorem~5 is that when $K$ is quadratic, $H$ is the
compositum of dihedral extensions of $\mathbb{Q}$. These are, in
principle, manageable as the splitting fields of polynomials with rational
coefficients.

Theorem~4 gives the number of 2-primary components of the group of ideal
classes of a quadratic field. For $\ell$ an odd prime, the question of
whether $\ell\mid h$ where $h$ is the class number is best handled by
examining the possibility of the existence of a cyclic unramified
extension $L/K$ of degree $\ell$. I shall treat this topic in §4 for the
case $\ell=3$, cf.\ Theorem~6.

In case $\mathfrak{H}\simeq\mathbb{Z}_{2}^{\,t-1}$ then
$\mathfrak{H}=\mathfrak{H}_{\Pi}$ and so $\mathfrak{G}=\mathfrak{H}$. The
absolute class field $H$ is then given by Theorem~1; all that is involved
is the simultaneous adjunction of square roots of rational integers. The
table below shows that for $|d|$ small these trivial cases are quite
common. (The first exception for $d>0$ is $d=145$,
$\mathfrak{H}\simeq\mathbb{Z}_{4}$.) One knows that there are only
finitely many quadratic fields with a given $\mathfrak{H}$ as class
group. I do not know which finite commutative groups are realizable as
the group of ideal classes of an imaginary quadratic field.

Let $K=\mathbb{Q}(\sqrt{d})$ be a quadratic field of discriminant
$d<0$. Below we list the structures of the groups of ideal classes for
all $d\geq -168$ and also $d=4d_{0}$ with $d_{0}>-100$ when these are
outside the list of $d\geq -168$ but give a group other than
$\mathbb{Z}_{2}^{\,t-1}$.

\begin{center}
\begin{tabular}{ll}
$\mathbb{Z}_{2}$:\ & $3,\ 4,\ 7,\ 8,\ 11,\ 19,\ 43,\ 67,\ 163,\ ?$\\
$\mathbb{Z}_{2}^{2}$:\ & $15,\ 20,\ 24,\ 35,\ 40,\ 51,\ 52,\ 88,\ 91,\ 123,\ 148,\dots$\\
$\mathbb{Z}_{2}^{3}$:\ & $84,\ 120,\ 132,\ 168,\dots$\\
& $420,\dots$\\
$\mathbb{Z}_{4}$:\ & $39,\ 55,\ 56,\ 68,\ 111,\ 136,\ 155,\dots$\\
$\mathbb{Z}_{8}$:\ & \\
$\mathbb{Z}_{2}\times\mathbb{Z}_{4}$:\ &\\
$\mathbb{Z}_{3}$:\ & $23,\ 31,\ 59,\ 107,\ 139,\dots$\\
$\mathbb{Z}_{2}\times\mathbb{Z}_{3}$:\ &\\
$\mathbb{Z}_{4}\times\mathbb{Z}_{3}$:\ &\\
$\mathbb{Z}_{5}$:\ &\\
$\mathbb{Z}_{2}\times\mathbb{Z}_{5}$:\ &\\
$\mathbb{Z}_{7}$:\ &\\
$\mathbb{Z}_{2}\times\mathbb{Z}_{7}$:\ &\\
$\mathbb{Z}_{11}$:\ &
\end{tabular}
\end{center}

% ---------------------------------------------------------------
\section{Unramified cubic extensions}\label{sec:VII-cubic}

Suppose $L/K$ is cyclic of degree $3$. Choose a generator $\sigma$ for
the Galois group of $L/K$ and let $\omega$ be a primitive cube root of
unity. We suppose $\omega\notin K$. Then $L(\omega)/K$ is cyclic of
degree $6$. We extend $\sigma$ to an automorphism of $L(\omega)/K$ by
putting $\omega^{\sigma}=\omega$. Let $\rho$ denote the unique
automorphism of $L(\omega)/K$ which has order $2$. There is an element
$\xi\in L(\omega)$ such that $\xi^{\sigma-1}=\omega$; $\xi$ is uniquely
determined by the choice of the pair $(\sigma,\omega)$ up to
multiplication by $\beta\in K(\omega)$.

Define quantities $a$, $e$, and $\theta$ by
\[
\theta=\xi+\xi^{\rho},\qquad
a=\xi^{1+\rho},\qquad
e=\xi^{2-\rho}+\xi^{2\rho-1}.
\]
It is easy to check that $a,e\in K$, $\theta\in L$. Indeed
\[
L=K(\theta)\qquad\text{where}\qquad\theta^{3}-3a\theta-ae=0.
\]

Now suppose $K/\mathbb{Q}$ is quadratic and $L/\mathbb{Q}$ is dihedral.
We can extend the non-trivial automorphism of $K/\mathbb{Q}$ to an
automorphism $\tau$ of $L(\omega)/\mathbb{Q}$ with the properties
\[
\omega^{\tau}=\omega^{2},\qquad\tau^{2}=1,\qquad
\rho\tau=\tau\rho,\qquad\sigma\tau=\tau\sigma^{2}.
\]
Put $\beta=\xi^{\tau-1}$. Then
\[
\beta^{\sigma}=\xi^{\tau\sigma-\sigma}
=\xi^{\sigma^{2}\tau}\xi^{-\sigma}
=(\omega^{2}\xi)^{\tau}(\omega^{-1}\xi)^{-1}
=(\omega\xi^{\tau})(\omega^{-1}\xi^{-1})
=\xi^{\tau-1}=\beta.
\]
Hence $\beta\in K(\omega)$, and we are free to replace $\xi$ by
$(1+\beta)\xi$ if $\beta\neq -1$; we replace $\xi$ by
$(\omega-\omega^{2})\xi$ if $\beta=-1$. With these replacements we may
assume that $\xi^{\tau}=\xi$ and $\xi$ is determined uniquely up to a
factor $\gamma\in K_{1}$ where $K=\mathbb{Q}(\sqrt{d})$,
$K_{1}=\mathbb{Q}(\sqrt{-3d})$. Moreover $a,e\in\mathbb{Q}$.

If $L/K$ is unramified then $L(\omega)/K(\omega)$ is unramified. Since
$\xi^{3}\in K(\omega)$, it follows that for a prime $\mathfrak{P}$ of
$K(\omega)$ if $\mathfrak{P}\mid\xi^{3}$ then
$\mathfrak{P}^{3}\mid\xi^{3}$. On the other hand $\xi^{3}\in K_{1}$;
hence if $\mathfrak{P}$ is a prime of $K_{1}$ and
$\mathfrak{P}\mid\xi^{3}$ then $\mathfrak{P}^{3}\mid\xi^{3}$. There is no
loss in generality in assuming that $\xi^{3}$ is an integer of $K_{1}$
which is not divisible by the cube of any integer (other than a unit) in
$K_{1}$. The norm from $K_{1}$ to $\mathbb{Q}$ of $\xi^{3}$ is
$\xi^{3(1+\rho)}=a^{3}$. The conclusion is that the rational prime $p$
divides $a\in\mathbb{Z}$ only if $p$ factors in $K_{1}$ into
non-principal prime ideals whose cubes are principal.

The above reasoning leads to

\begin{namedstmt}{THEOREM 6}
Let $K=\mathbb{Q}(\sqrt{d})$ be an imaginary quadratic field with class
number $h$. Put $K_{1}=\mathbb{Q}(\sqrt{-3d})$. Let $h_{1}$ be the class
number of $K_{1}$, $\varepsilon$ the fundamental unit of $K_{1}$, $a$ the
norm of $\varepsilon$, and $e$ the trace of $\varepsilon$. Then $3\mid h$
iff at least one of the following conditions holds
\[
\begin{aligned}
&(\mathrm{i}_{1})\ a=+1,\ e\equiv\pm 2\bmod 27\\
&(\mathrm{i}_{2})\ e\equiv 0\bmod 9\\
&(\mathrm{i}_{3})\ a=-1,\ e\equiv\pm 4\bmod 9\\
&(\mathrm{ii})\ 3\mid h_{1}.
\end{aligned}
\]
If (i) holds then $K(\theta)/K$ is an unramified cyclic cubic extension
of $K$ where $\theta^{3}-3a\theta-ae=0$.
\end{namedstmt}

\noindent\emph{Remark.} Cases $(\mathrm{i}_{2})$ and
$(\mathrm{i}_{3})$ can only occur if $3\mid d$.

\begin{proof}[Proof.]
If $3\mid h$ then $K$ has an unramified cyclic cubic extension
$L=K(\theta)$ where $\theta=\xi+\xi^{\rho}$ and $\xi^{3}\in K_{1}$. As we
saw above, either (ii) holds or $\xi^{3}$ is a unit. If $\xi^{3}$ is a
unit then, modifying $\xi$ by a factor $\pm\varepsilon^{k}$, we can
suppose that $\xi^{3}=\varepsilon$ or $\varepsilon^{\rho}$. The extension
$K(\theta)/K$ where $\theta^{3}-3a\theta-ae=0$ can only ramify at primes
dividing $3$ in $K$. If $3\nmid d$ then $3$ does not ramify in
$K/\mathbb{Q}$, and hence $K(\theta)/K$ is unramified iff
$\mathbb{Q}(\theta)/\mathbb{Q}$ is unramified at $3$. In turn
$\mathbb{Q}(\theta)/\mathbb{Q}$ is unramified at $3$ iff
$e\equiv\pm 2\bmod 27$ (because when $3\nmid d$ we must have $a=1$ and
$3\nmid e$). When $3\mid d$ the analysis is more complicated, but in all
cases condition (i) is necessary and sufficient that $K(\theta)/K$ be
unramified. Finally, if $3\mid h_{1}$ then $K_{1}$ has an unramified
cyclic cubic extension, but a real quadratic field
$\mathbb{Q}(\sqrt{d_{1}})$ can have such an extension only if the class
number of the imaginary quadratic field $\mathbb{Q}(\sqrt{-3d_{1}})$ is
divisible by $3$ since $\mathbb{Q}(\sqrt{-3d_{1}})$ has no units other
than $\pm 1$.
\end{proof}

\subsection*{Examples.}

\noindent
($\mathrm{i}_{1}$) This is the case most frequently encountered, e.g.\
$d=-23$, $e=25$; $d=-31$, $e=29$; $d=-116$, $e=56$; $d=-152$,
$e=2050$; $d=-1979$, $e=79$. The case $d=-107$ is more interesting. One
has $N\bigl(\frac{11+9\sqrt{-107}}{2}\bigr)=13^{3}$ and $h=3$. Also
$h_{1}=3$, $\varepsilon=215+12\sqrt{321}$. Using ($\mathrm{i}_{1}$) we
get $H=\mathbb{Q}(\sqrt{-107},\theta)$ where $\theta^{3}-3\theta-430=0$.
Note that $N\bigl(\frac{17+\sqrt{321}}{2}\bigr)=(-2)^{3}$ but
$\xi^{3}=(17+\sqrt{321})/2$ does not yield an unramified extension since
$\xi^{6}\not\equiv 1\bmod\mathfrak{P}_{3}^{3}$ where $\mathfrak{P}_{3}$
is the prime dividing $3$ in $K_{1}$. The absolute class field of $K_{1}$
is $H_{1}=\mathbb{Q}(\sqrt{321},\theta_{1})$ where
$\theta_{1}^{3}-3\cdot 13\,\theta_{1}-11=0$.

\medskip\noindent
($\mathrm{i}_{2}$) $d=-3\cdot 7\cdot 11$. Here
$\mathfrak{H}\simeq\mathbb{Z}_{2}^{2}\times\mathbb{Z}_{3}$. We have
$h_{1}=1$ and $\varepsilon=(9+\sqrt{77})/2$. The absolute class field is
$H=\mathbb{Q}(\sqrt{-3},\sqrt{-7},\sqrt{-11},\varepsilon^{1/3})$.

\medskip\noindent
($\mathrm{i}_{3}$) $d=-3\cdot 29$. Here
$\mathfrak{H}\simeq\mathbb{Z}_{2}\times\mathbb{Z}_{3}$, $h_{1}=1$,
$\varepsilon=(5+\sqrt{29})/2$. The absolute class field is
$H=\mathbb{Q}(\sqrt{-3},\sqrt{29},\varepsilon^{1/3})$.

\medskip\noindent
(ii) $d=-3\cdot 229$. One has $N(26+3\sqrt{-687})=19^{3}$, hence
$3\mid h$. Also $2\mid h$ and, in fact,
$\mathfrak{H}\simeq\mathbb{Z}_{2}\times\mathbb{Z}_{3}$. We have
$K_{1}=\mathbb{Q}(\sqrt{229})$, $\varepsilon=(15+\sqrt{229})/2$, so none
of the cases (i) apply. On the other hand $N(6+\varepsilon)=5^{3}$ and
$(6+\varepsilon)^{2}\equiv 1\bmod 9$. The absolute class field to $K$ is
$H=\mathbb{Q}(\sqrt{-3},\sqrt{229},(6+\varepsilon)^{1/3})$. At the same
time, since $h_{1}=3$, the absolute class field to $K_{1}$ is
$H_{1}=\mathbb{Q}(\sqrt{229},\theta_{1})$ where
$\theta_{1}^{3}-3\cdot 19\,\theta_{1}-52=0$.

% ---------------------------------------------------------------
\section{Unramified cyclic extensions of degree 4}\label{sec:VII-quartic}

Let $K=\mathbb{Q}(\sqrt{d})$ be a quadratic field of discriminant $d$.
Suppose $L/K$ is an unramified cyclic extension of degree $4$. We know
that the Galois group of $L/\mathbb{Q}$ is given by generators $\sigma$
and $\tau$, where $K$ is the fixed field of $\sigma$, satisfying the
relations $\sigma^{4}=1=\tau^{2}$, $\sigma\tau=\tau\sigma^{3}$. Let $F$
be the fixed field of $\sigma^{2}$. Then $F$ is an unramified quadratic
extension of $K$; hence $F\subset G$ where $G$ is the genus field. It
follows from Theorem~1 that
$F=\mathbb{Q}(\sqrt{d_{+}},\sqrt{d_{-}})$ where $d=d_{+}d_{-}$; the
factors are distinguished by
\[
\sqrt{d_{+}}{}^{\,\tau}=+\sqrt{d_{+}},\qquad
\sqrt{d_{-}}{}^{\,\tau}=-\sqrt{d_{-}}.
\]
(If $d<0$ this means $d_{+}>0$ and $d_{-}<0$.) Take $R$ the fixed field
of $\tau$. Then $R$ is quadratic over $\mathbb{Q}(\sqrt{d_{+}})$; so we
can write $R=\mathbb{Q}(\sqrt{d_{+}},\sqrt{\mu})$ where
$\mu\in\mathbb{Q}(\sqrt{d_{+}})$ and
$\sqrt{\mu}^{\,\sigma^{2}}=-\sqrt{\mu}$. On the other hand $R/\mathbb{Q}$
is not normal; hence $\mu\notin\mathbb{Q}$ and
$L=\mathbb{Q}(\sqrt{d_{-}},\sqrt{\mu})$. One finds easily that
$\sqrt{\mu}^{\,1+\sigma}$ reverses sign under the actions of both
$\sigma$ and $\tau$. It follows that
$\sqrt{\mu}^{\,1+\sigma}=m\sqrt{d_{-}}$ where $m\in\mathbb{Q}$. The
result is
\[
N\mu=\mu^{1+\sigma}=m^{2}d_{-}
\]
where $N$ denotes the norm from $\mathbb{Q}(\sqrt{d_{+}})$ to
$\mathbb{Q}$.

There is obviously no loss in generality in assuming that $\mu$ is a
square-free integer of $\mathbb{Q}(\sqrt{d_{+}})$. This fixes $\mu$ up to
multiplication by the square of a unit if $\mathbb{Q}(\sqrt{d_{+}})$ has
all its ideals principal. In general, what one can say is given in

\begin{namedstmt}{LEMMA 3}
Let $p$ be a rational prime; if $p\mid d_{-}$ then $p$ splits in
$\mathbb{Q}(\sqrt{d_{+}})$. Write $d_{-}=d_{0}$ or $d_{-}=4d_{0}$,
where $4\nmid d_{0}$. The factorization of $\mu$ in
$\mathbb{Q}(\sqrt{d_{+}})$ is $(\mu)=\underline{\mathfrak{m}}^{2}
\underline{\mathfrak{d}}$ where $N\underline{\mathfrak{d}}=(d_{0})$ and
$\underline{\mathfrak{m}}$ has no principal ideal factors.
\end{namedstmt}

\begin{proof}[Proof.]
Let $\underline{\mathfrak{p}}$ be a prime of $\mathbb{Q}(\sqrt{d_{+}})$
and suppose $\underline{\mathfrak{p}}^{\,k}$ exactly divides $\mu$. In
order that $F(\sqrt{\mu})$ be unramified over $F$ it is necessary that
$\underline{\mathfrak{p}}^{\,k}$ be a square in $F$. Hence $k$ must be
even unless $\underline{\mathfrak{p}}\mid d_{-}$. Suppose
$\underline{\mathfrak{p}}\mid d_{0}$ and put $p^{f}=N\underline{
\mathfrak{p}}$. Taking norms, we have that $p^{kf}\mid m^{2}d_{-}$. Thus
$kf$ must be odd. Therefore $f=1$ ($p$ splits) and $k$ is odd. Everything
is proved except for the case $d_{-}=4d_{0}$, $d_{0}$ odd, and
$\underline{\mathfrak{p}}\mid 2$. This requires a detailed examination of
the equation $x^{2}-\mu=0$ locally at the primes dividing $2$ in $F$.
The desired result is obtained, but we omit the details.
\end{proof}

A simple description of the whole story is

\begin{namedstmt}{THEOREM 7}
Let $K=\mathbb{Q}(\sqrt{d})$ be a quadratic field of discriminant $d$. A
necessary and sufficient condition that $K$ have an unramified cyclic
extension of degree $4$ is that there exist an admissible factorization
$d=d_{+}d_{-}$ such that $d_{-}$ is a norm from
$\mathbb{Q}(\sqrt{d_{+}})$ and $2$ splits in
$\mathbb{Q}(\sqrt{d_{+}})$ if $2\mid d_{-}$.
\end{namedstmt}

\begin{proof}[Proof.]
If such an extension exists we have $d_{-}=N(m^{-1}\mu)$ and
Lemma~3 says that $2$ splits in $\mathbb{Q}(\sqrt{d_{+}})$ if
$2\mid d_{-}$. Now suppose the condition is fulfilled. Then
$d_{0}=\alpha^{1+\sigma}$ for some $\alpha\in\mathbb{Q}(\sqrt{d_{+}})$.
The prime factorization of $\alpha$ must have the form
$(\alpha)=\underline{\mathfrak{d}}\,\underline{\mathfrak{m}}^{1-\sigma}$.
Put $m_{0}=N\underline{\mathfrak{m}}$; then $\mu=m_{0}\alpha$ has all the
properties previously given (if this $\mu$ is not square-free, it can be
made so by modifying $\underline{\mathfrak{m}}$). We may also assume that
$\underline{\mathfrak{m}}$ is relatively prime to $2$, for suppose
$\underline{\mathfrak{m}}$ is exactly divisible by
$\underline{\mathfrak{p}}^{k}$ where $\underline{\mathfrak{p}}\mid 2$.
Since $\mu$ is square-free, $\underline{\mathfrak{p}}$ is not principal.
Therefore $\underline{\mathfrak{p}}^{1+\sigma}=(2)$ and there is an
integer $\beta\in\mathbb{Q}(\sqrt{d_{+}})$ such that
$(\beta)=\underline{\mathfrak{p}}^{\sigma}\underline{\mathfrak{q}}$
where $\underline{\mathfrak{q}}$ is relatively prime to $2$. We replace
$\underline{\mathfrak{m}}$ by $(\beta/2)^{k}\underline{\mathfrak{m}}$,
i.e.\ we change $\mu$ to $(\beta/2)^{2k}\mu$, with no loss. The equation
$x^{2}-\mu=0$ does not ramify over $F$ at any prime not dividing $2$
since $(\mu)=\underline{\mathfrak{m}}^{2}\underline{\mathfrak{d}}$ is the
square of an ideal in $F$. The ramification at $2$ is handled this way:
we have
\[
\mu=(a+b\sqrt{d_{+}})/2,\qquad a,b\in\mathbb{Z}\ \text{both even or both odd},
\]
and the calculation of $N\mu$ gives
\[
4m_{0}^{2}d_{0}=a^{2}-b^{2}d_{+};\qquad(m_{0}\text{ is odd}).
\]
Let $b>0$ be fixed such that $(\pm a,b)$ is a solution to the norm
equation. We can always choose the sign of $a$ so that $x^{2}-\mu=0$ does
not ramify over $F$ at primes dividing $2$. The proof requires an
analysis of cases. Case I is $d_{+}\equiv 1\bmod 4$, $d_{-}\equiv
1\bmod 4$; here the condition is that $\mu$ be a quadratic residue mod
$4$ in $\mathbb{Q}(\sqrt{d_{+}})$. The remaining cases are
\[
\begin{aligned}
\text{II:}&\quad d_{+}\equiv 8\bmod 32,\quad d_{-}\equiv 1\bmod 4\\
\text{III:}&\quad d_{+}\equiv 1\bmod 8,\quad d_{-}\equiv -4\bmod 16\\
\text{III$'$:}&\quad d_{+}\equiv 1\bmod 8,\quad d_{-}\equiv -8\bmod 32.
\end{aligned}
\]
The ramification conditions become complicated when $2\mid d_{-}$, and we
omit the tedious details of the case analyses.

The unramified cyclic extension of degree 4 is of the form
$L=\mathbb{Q}(\sqrt{d_{-}},\sqrt{\mu})$.
\end{proof}

Here are some examples, including all known cases where $d<0$, $4\mid d$,
and $\mathfrak{H}\simeq\mathbb{Z}_{4}$:

\begin{center}
\begin{tabular}{rrrl}
\multicolumn{1}{c}{$d_{+}$} & \multicolumn{1}{c}{$d_{-}$} & 
\multicolumn{1}{c}{$\mu$} & \multicolumn{1}{c}{$\mathfrak{H}$}\\
\hline
$-56$ & $8$ & $-7$:\ \ $-1+2\sqrt{2}$ & $\mathbb{Z}_{4}$\\
$-184$ & $8$ & $-23$:\ \ $-3+4\sqrt{2}$ & $\mathbb{Z}_{4}$\\
$-248$ & $8$ & $-31$:\ \ $1+4\sqrt{2}$ & $\mathbb{Z}_{8}$\\
$-568$ & $8$ & $-71$:\ \ $-1+6\sqrt{2}$ & $\mathbb{Z}_{4}$\\
$-68$ & $17$ & $-4$:\ \ $4+\sqrt{17}$ & $\mathbb{Z}_{4}$\\
$-64$ & $41$ & $-4$:\ \ $32+5\sqrt{41}$ & $\mathbb{Z}_{8}$\\
$-260$ & $65$ & $-4$:\ \ $8+\sqrt{65}$ & $\mathbb{Z}_{2}\times\mathbb{Z}_{4}$\\
$-292$ & $73$ & $-4$:\ \ $1068+125\sqrt{73}$ & $\mathbb{Z}_{4}$\\
$-356$ & $89$ & $-4$:\ \ $500+53\sqrt{89}$ & $\mathbb{Z}_{4}$\\
$-772$ & $193$ & $-4$:\ \ $1764132+126985\sqrt{193}$ & $\mathbb{Z}_{4}$\\
$-136$ & $17$ & $-8$:\ \ $(3+\sqrt{17})/2$ & $\mathbb{Z}_{4}$\\
$-328$ & $41$ & $-8$:\ \ $(19+3\sqrt{41})/2$ & $\mathbb{Z}_{4}$\\
$-55$ & $5$ & $-11$:\ \ $3+2\sqrt{5}$ & $\mathbb{Z}_{4}$\\
$-95$ & $5$ & $-19$:\ \ $-1+2\sqrt{5}$ & $\mathbb{Z}_{8}$\\
$-155$ & $5$ & $-31$:\ \ $-7+4\sqrt{5}$ & $\mathbb{Z}_{4}$\\
$-39$ & $13$ & $-3$:\ \ $(-1+\sqrt{13})/2$ & $\mathbb{Z}_{4}$\\
$-111$ & $37$ & $-3$:\ \ $(-5+\sqrt{37})/2$ & $\mathbb{Z}_{4}$\\
$-145$ & $29$ & $-5$:\ \ $11+2\sqrt{29}$ & $\mathbb{Z}_{4}$\\
\end{tabular}
\end{center}

(In the real case one distinguishes $d_{+}$ from $d_{-}$ by testing: here
$7$ remains prime in $\mathbb{Q}(\sqrt{145})$ and $\mathbb{Q}(\sqrt{5})$
but splits in $\mathbb{Q}(\sqrt{29})$; hence $\sqrt{29}$ is fixed under
$\tau$; cf.\ §\ref{sec:VII-genus}.)
